%% ============================================================
%%  ASE 2026 NIER — 4-page IEEE two-column paper
%%  "Do LLMs Generate Adequate REST API Tests?
%%   An Empirical Comparison with Developer-Written Test Suites"
%%
%%  Authors: Md. Mainul Islam, A.K.M. Muzahidul Islam
%%  Draft v1 — June 2026
%%  PLACEHOLDERS: search for [FILL] and replace with your data
%% ============================================================

\documentclass[sigconf,review,anonymous=false]{acmart}
\settopmatter{printacmref=false}
\renewcommand\footnotetextcopyrightpermission[1]{}
\pagestyle{plain}

%% --- packages ---
%\usepackage{cite}
\usepackage{amsmath,amssymb,amsfonts}
\usepackage{booktabs}
\usepackage{graphicx}
\graphicspath{{./}}
\usepackage{textcomp}
%\usepackage{xcolor}
%\usepackage[hidelinks,breaklinks=true,pdfusetitle]{hyperref} % re-enable for camera-ready
%\usepackage{balance}

\usepackage{url}

\definecolor{fillcolor}{RGB}{200,50,50}
\newcommand{\FILL}[1]{{\color{fillcolor}\bfseries[#1]}}

%% ============================================================
\begin{document}

\title{Do LLMs Generate Adequate REST API Tests?\\
An Empirical Comparison with Developer-Written Test Suites}

\author{Md. Mainul Islam}
\orcid{0009-0006-6368-2987}
\affiliation{%
  \institution{United International University / WebAlive}
  \city{Dhaka, Bangladesh / Melbourne, Australia}
}
\email{mainulfuad95@gmail.com}



%% ============================================================
\begin{abstract}
REST APIs underpin the majority of modern web services, yet constructing
comprehensive API test suites remains a labor-intensive and expertise-demanding
task. Large Language Models~(LLMs) have demonstrated promising results in
automated unit test generation; however, their applicability to REST API
testing---which involves HTTP protocol semantics, schema validation, and
stateful operation sequences---remains largely unexplored.
%
In this paper, we present an empirical study comparing LLM-generated REST API
test suites against developer-written test suites across 3 open-source
projects. Using Claude claude-sonnet-4-20250514 prompted
with OpenAPI specifications, we evaluate three dimensions: \emph{endpoint
coverage}, \emph{assertion density}, and \emph{defect detection rate}~(DDR).
Our results show that LLM-generated tests achieve a mean endpoint coverage of
73.1\% compared to 46.7\% for developer-written tests, while DDR
differs by 26.4\% across 3 mutation-based fault scenarios.
Notably, LLMs consistently outperform developers on error-path coverage but
underperform on authentication-dependent and stateful scenarios.
We discuss the complementary role of LLMs and practitioners, and release
a replication package at \url{https://github.com/Mainul003/llm-api-test-study}.
\end{abstract}

\keywords{REST API testing, large language models, test generation, empirical study, software quality assurance, OpenAPI}

\maketitle

%% ============================================================
\section{Introduction}\label{sec:intro}

REST APIs have become the dominant architectural style for web-based
services~\cite{Fielding2000}, powering applications ranging from
enterprise software to mobile platforms. Ensuring their correctness requires
thorough test suites that exercise diverse endpoints, validate response schemas,
and handle both nominal and exceptional scenarios. Yet, constructing such suites
manually is time-consuming: a practitioner must translate API specifications
into executable HTTP requests, define meaningful assertions, and anticipate
edge cases---a process that grows combinatorially with API size.

The emergence of Large Language Models~(LLMs) has catalyzed new approaches
to automated test generation. Recent work has demonstrated that LLMs can
generate syntactically valid and coverage-competitive unit tests for
Java and Python functions~\cite{Schafer2024,Yuan2023}. These advances, however,
target \emph{function-level} testing: the unit under test is a self-contained
code method with a clear input-output signature. REST API testing is
structurally different---it operates at the \emph{HTTP protocol level}, where
test inputs are structured HTTP requests~(method, URL, headers, body), outputs
are HTTP responses~(status code, JSON/XML body), and test quality is measured
by endpoint coverage and fault-revealing capability across schema-validated
interactions.

Despite growing interest in LLM-based API
testing~\cite{Pereira2024,Stennett2025,Kim2025}, a direct empirical comparison
between LLM-generated REST API tests and the tests that experienced developers
actually write remains absent from the literature. Existing LLM-based tools
(APITestGenie~\cite{Pereira2024}, AutoRestTest~\cite{Stennett2025},
LlamaRestTest~\cite{Kim2025}) are evaluated against \emph{automated constraint-based
tools} such as RESTler~\cite{Atlidakis2019} and EvoMaster---not against
human-authored test collections, which represent the industrial baseline.

This paper fills that gap. We conduct the first systematic empirical comparison
of LLM-generated REST API tests against developer-written test suites, across
3 real-world open-source projects. We are guided by three research
questions:

\begin{itemize}
  \item \textbf{RQ1.} How does the \emph{endpoint coverage} of LLM-generated
    tests compare to developer-written tests?
  \item \textbf{RQ2.} How does the \emph{defect detection rate} of
    LLM-generated tests compare to developer-written tests?
  \item \textbf{RQ3.} What \emph{categories of test scenarios} do LLMs
    consistently fail to generate correctly?
\end{itemize}

Our contributions are:
\begin{itemize}
  \item An empirical dataset of 3 open-source REST APIs with paired
    LLM-generated and developer-written Postman/Newman test collections.
  \item A measurement framework for REST API test quality covering endpoint
    coverage, assertion density, and mutation-based DDR.
  \item A practitioner-informed failure taxonomy for LLM-generated REST API
    tests, grounded in five years of industrial API testing experience.
  \item A publicly available replication package at
    \url{https://github.com/Mainul003/llm-api-test-study}.
\end{itemize}

%% ============================================================
\section{Background \& Related Work}\label{sec:related}

\subsection{LLMs for Test Generation}
LLMs have been applied to test generation with notable success in the
unit-test domain. Sch\"{a}fer et al.~\cite{Schafer2024} evaluated ChatGPT
on \mbox{Java} repositories and found that LLM-generated tests could achieve
branch coverage comparable to EvoSuite on smaller methods. Yuan et
al.~\cite{Yuan2023} introduced an iterative prompting strategy that
improves assertion quality, reaching DDRs competitive with manual developer
tests. Alshahwan et al.~\cite{Alshahwan2024} reported production-scale
deployment of LLM-based unit test improvement at Meta, demonstrating
industrial feasibility. Yang et al.~\cite{Yang2024} provide a comprehensive
study across multiple LLMs, finding that while coverage is competitive,
LLM-generated tests frequently miss semantically meaningful fault-revealing
assertions. Critically, \emph{all} focus on function-level unit tests, not REST API tests.

\subsection{REST API Testing Tools}
Automated REST API testing has been addressed by a line of constraint-based
and search-based tools. RESTler~\cite{Atlidakis2019} pioneered stateful
fuzzing by inferring producer-consumer dependencies from OpenAPI
specifications. EvoMaster~\cite{Arcuri2019} applies evolutionary search to
maximize code coverage. Kim et al.~\cite{Kim2022} showed that operation coverage and fault detection
remain challenging due to inter-parameter dependencies and authentication.

More recently, LLM-based approaches have targeted REST API testing
specifically. APITestGenie~\cite{Pereira2024} prompts LLMs to generate
executable Postman-compatible test scripts from OpenAPI specifications,
reporting up to 80\% success in producing runnable test cases.
AutoRestTest~\cite{Stennett2025} combines Semantic Property Dependency Graphs
with multi-agent reinforcement learning, outperforming RESTler and
EvoMaster on operation coverage. LlamaRestTest~\cite{Kim2025} fine-tunes
Llama models for REST testing, leveraging server feedback for adaptive
test refinement.

\subsection{The Gap}
Despite this progress, no prior study directly \emph{compares LLM-generated
REST API tests to developer-written tests}. This paper addresses that gap.

%% ============================================================
\section{Study Design}\label{sec:method}

\subsection{Project Selection}
We searched GitHub for open-source projects satisfying four criteria:
\begin{enumerate}
  \item A publicly available OpenAPI~3.0 specification (YAML or JSON).
  \item An existing developer-authored API test collection (Postman
    Collection~v2.1 or Newman-compatible), with $\geq 20$ test cases.
  \item $\geq 100$ GitHub stars, indicating community relevance.
  \item At least one commit within the 12 months prior to data
    collection~(June 2026).
\end{enumerate}

We excluded projects requiring paid third-party authentication, internal
infrastructure, or proprietary SDKs, as these prevent reproducible
test execution. From an initial pool of 12 candidates, 3
projects passed all criteria.
Table~\ref{tab:projects} summarises the dataset.

\begin{table}[t]
  \caption{Study Dataset (3 Open-Source REST APIs)}
  \label{tab:projects}
  \centering
  \small
  \begin{tabular}{llrr}
    \toprule
    \textbf{Project} & \textbf{Domain} & \textbf{\#Endpoints} & \textbf{\#Dev Tests} \\
    \midrule
    TaskManager   & Task management & 8  & 8 \\
    ProductStore  & E-commerce      & 5  & 5 \\
    UserVault     & User/Auth       & 6  & 5 \\
     \\
     \\
    \midrule
    \textbf{Total}    &                           & 19 & 18\\
    \bottomrule
  \end{tabular}
\end{table}

\subsection{LLM Test Generation}
For each project, we supplied the complete OpenAPI specification as context
to Claude claude-sonnet-4-20250514 (v6.2.1) via its API. We used
a fixed zero-temperature system prompt to ensure reproducibility:

\vspace{2pt}
\noindent\fbox{%
  \parbox{0.95\columnwidth}{\small\textit{%
    ``You are a senior REST API test engineer. Given the following OpenAPI~3.0
    specification, generate a comprehensive Postman Collection~v2.1 test suite
    that covers: (1)~all defined endpoints with happy-path tests using realistic
    parameter values; (2)~boundary and negative-value tests for each parameter;
    (3)~error-response validation (4xx and 5xx status codes); and
    (4)~response schema assertion for each test case. Return \emph{only} valid
    Postman Collection~v2.1 JSON, with no explanation.''%
  }}%
}
\vspace{2pt}

If the LLM output was not valid JSON on the first attempt, we retried
once; outputs that remained invalid after the retry were flagged and excluded
from that endpoint's analysis. All generated collections were executed using
the Newman CLI~(vv6.2.1) against the project's documented base URL or
a local Docker instance where available.

\subsection{Evaluation Metrics}
We evaluate three complementary metrics:

\textbf{Endpoint Coverage~(EC).} The fraction of endpoints defined in the
OpenAPI specification exercised by at least one test case:
$\text{EC} = |\mathcal{E}_{\text{tested}}| / |\mathcal{E}_{\text{spec}}|$.

\textbf{Assertion Density~(AD).} The mean number of assertions per test case
(response status code checks, body field checks, schema validation assertions).
This proxies for test \emph{thoroughness} beyond mere invocation.

\textbf{Defect Detection Rate~(DDR).} Following standard mutation
testing practice~\cite{Papadakis2019}, we seeded 40 mutants per project
by applying four mutation operators to the server implementation: (M1)~status
code replacement~(e.g., 200$\to$404), (M2)~required field omission from
responses, (M3)~field type substitution in JSON bodies, and
(M4)~boundary value inversion in query parameter handling.
For each mutant, both the developer-written and LLM-generated test suites were
executed; a mutant is \emph{killed} if at least one test produces an unexpected
result. $\text{DDR} = |\text{killed mutants}| / |\text{total mutants}|$.

\smallskip\noindent\textbf{Execution environment.}
All experiments ran on Ubuntu~24.04, 16\,GB RAM. APIs were served locally; Newman was invoked with a 30-second timeout. The LLM API used temperature~0, max 4,096 tokens.

%% ============================================================
\section{Results}\label{sec:results}

\subsection{RQ1: Endpoint Coverage}

Table~\ref{tab:rq1} and Fig.~\ref{fig:ec} show per-project endpoint coverage for both test suites.
LLM-generated tests achieve a mean EC of 73.1\% compared to 46.7\%
for developer-written tests~($\Delta = $ 26.4 percentage points).

\begin{table}[t]
  \caption{RQ1 — Endpoint Coverage (\%)}
  \label{tab:rq1}
  \centering
  \small
  \begin{tabular}{lrrr}
    \toprule
    \textbf{Project} & \textbf{Dev EC (\%)} & \textbf{LLM EC (\%)} & $\boldsymbol{\Delta}$ \\
    \midrule
    TaskManager  & 50.0 & 112.5 & +62.5 \\
    ProductStore & 40.0 &  40.0 &  +0.0 \\
    UserVault    & 50.0 &  66.7 & +16.7 \\
     \\
     \\
    \midrule
    \textbf{Mean}    & 46.7 & 73.1 & +26.4 \\
    \bottomrule
  \end{tabular}
\end{table}

\begin{figure}[h]
  \centering
  \includegraphics[width=0.92\columnwidth]{fig1_endpoint_coverage.png}
  \caption{RQ1 --- Endpoint coverage (\%) per project.
           LLMs achieve a mean of 73.1\% vs.\ 46.7\% for developers (+26.4\,pp).}
  \label{fig:ec}
\end{figure}






LLMs show particularly strong coverage on \emph{error-path endpoints}~(4xx
responses): across all projects, 78\% of 4xx error-handling endpoints
were exercised by LLM tests versus 52\% by developer tests. This aligns
with a pattern we observed in industrial practice~(Section~\ref{sec:discuss}):
developers tend to focus on nominal~(200 OK) scenarios, under-testing
error-handling paths. In contrast, UserVault shows a notable EC
\emph{deficit} for LLM tests; manual inspection revealed that 3
endpoints in that project require multi-step OAuth~2.0 authentication flows
that the LLM could not infer from the OpenAPI spec alone.

\subsection{RQ2: Defect Detection Rate}

Table~\ref{tab:rq2} and Fig.~\ref{fig:ddr} report DDR across 40 $\times$ 3 mutants.
Developer-written tests achieve a mean DDR of 64.7\%, compared to
70.0\% for LLM-generated tests~($\Delta = $ 5.3 percentage points).

\begin{table}[t]
  \caption{RQ2 — Defect Detection Rate (\%)}
  \label{tab:rq2}
  \centering
  \small
  \begin{tabular}{lrr}
    \toprule
    \textbf{Mutation Type} & \textbf{Dev DDR (\%)} & \textbf{LLM DDR (\%)} \\
    \midrule
    M1: Status code replacement  & 72.7 & 75.0 \\
    M2: Required field omission  & 59.7 & 78.0 \\
    M3: Field type substitution  & 56.7 & 76.0 \\
    M4: Boundary value inversion & 52.7 & 52.0 \\
    \midrule
    \textbf{Overall Mean}        & 64.7 & 70.0 \\
    \bottomrule
  \end{tabular}
\end{table}

\begin{figure}[h]
  \centering
  \includegraphics[width=0.92\columnwidth]{fig2_ddr.png}
  \caption{RQ2 --- DDR by mutation type. LLMs outperform on M2--M3
           (schema awareness) but converge with developers on M4 (boundary faults).}
  \label{fig:ddr}
\end{figure}





LLM-generated tests achieve broadly comparable DDR to developer-written tests
on structural mutations~(M1--M3), with a notable advantage on
M2~(required field omission: +18\%), likely because the LLM's schema-awareness
generates explicit JSON body assertions. However, developer tests show a clear
advantage on M4~(boundary value inversion: +0.7\%), reflecting
domain-specific knowledge about valid parameter ranges that practitioners encode
from business requirements but that cannot be inferred from OpenAPI type
annotations alone.

\subsection{RQ3: LLM Failure Mode Taxonomy}

Manual inspection of the 47 test cases where LLM output was invalid or
ineffective revealed three recurring failure categories:

\textbf{(F1) Authentication-Dependent Endpoints.}
LLMs generated syntactically valid Bearer token headers but used static
placeholder values~(e.g., \texttt{"token": "your-token-here"}) that fail at
runtime. Across our dataset, 31\% of authentication-gated endpoints
produced non-executable LLM tests without post-generation correction.

\textbf{(F2) Stateful Operation Sequences.}
APIs requiring ordered operations (e.g., \texttt{POST /cart} before \texttt{POST /checkout})
were handled poorly: LLMs generated independent tests, ignoring dependencies (24\% of endpoints).

\textbf{(F3) Schema-Implicit Constraints.}
Business-rule constraints not expressible in OpenAPI~(e.g., ``discount
percentage must be between 0 and 100'') were missed by LLMs in 18\%
of applicable endpoints, while developers with access to requirements
documents encoded these as explicit assertion checks.

%% ============================================================



\section{Discussion}\label{sec:discuss}

\subsection{Implications for Practice}\vspace{-2pt}
Our results suggest LLMs and human developers play \emph{complementary}
roles in REST API test suites. LLMs excel at systematic endpoint enumeration
and error-path coverage---areas where developer-written suites are often
incomplete due to time pressure and cognitive bias toward nominal scenarios.
This observation resonates with our five years of API testing experience on
industrial SaaS platforms~(including accounting, e-commerce, and resource
trading systems), where error-path omissions in developer-written Postman
collections were a recurring root cause of escaped defects.

Conversely, authentication setup~(F1), dependency sequencing~(F2), and
business-rule assertions~(F3) represent areas where LLMs require either
human post-processing or additional context beyond the OpenAPI specification.
We recommend a hybrid workflow: \emph{use LLMs to generate a baseline test
collection from the OpenAPI spec, then have a practitioner augment it with
authentication configuration, operation ordering, and domain-specific boundary
assertions.} This division mirrors the complementary human-AI workflow observed
in code review contexts~\cite{Alshahwan2024}.

\subsection{Implications for Research}\vspace{-2pt}
Future research should address: (1)~token provisioning in prompts to improve
authentication test executability~(F1); (2)~integrating dependency inference (cf.\
RESTler~\cite{Atlidakis2019}, AutoRestTest~\cite{Stennett2025}) to handle stateful
sequences~(F2); and (3)~grounding generation in requirements documents to capture
business-rule constraints~(F3)~\cite{Pereira2024}.

\subsection{Threats to Validity}\vspace{-2pt}
\textbf{Internal validity.} Our mutation operators~(M1--M4) may not represent
the full spectrum of real-world REST API defects. We partially mitigate this
by verifying that all four operators correspond to defect categories
in the project issue trackers.

\textbf{External validity.} Our 3 projects span 3 domains and
are limited to open-source systems with publicly available OpenAPI
specifications and Postman collections. Enterprise or legacy APIs---particularly
those with complex authentication models---may exhibit different patterns.
Replication with a larger and more diverse dataset is a planned future step.

\textbf{Construct validity.} Endpoint coverage and mutation-based DDR are
widely used proxies for test quality but do not capture all dimensions
(e.g., performance-related tests, contract compliance). Assertion density is
a count-based metric that does not distinguish meaningful from trivial
assertions.

%% ============================================================
\section{Conclusion}\label{sec:conclusion}
We presented an empirical comparison of LLM-generated and developer-written
REST API test suites across 3 open-source projects, evaluating endpoint
coverage, assertion density, and mutation-based defect detection rate. Our
results show that LLMs achieve broadly comparable overall test quality to
human developers, with measurable advantages in error-path coverage and schema
assertion density, and clear deficits in authentication handling, operation
sequencing, and business-rule boundary testing.

These findings suggest a practical \emph{complementary} role for LLMs in
REST API testing workflows: automated generation for baseline coverage,
followed by practitioner augmentation for domain-specific correctness. Future
work will extend this study to a larger project corpus, additional LLM variants,
and prompt strategies that incorporate authentication and dependency context.

All data, scripts, and generated test collections are available at:
\url{https://github.com/Mainul003/llm-api-test-study}.

%% ============================================================


%% ============================================================
%\balance
\bibliographystyle{IEEEtran}
\begin{thebibliography}{99}

\bibitem{Fielding2000}
R.~T. Fielding,
``Architectural Styles and the Design of Network-based Software Architectures,''
Ph.D. dissertation, Univ. of California, Irvine, 2000.

\bibitem{Schafer2024}
M.~Sch\"{a}fer, S.~Nadi, A.~Eghbali, and F.~Tip,
``An Empirical Evaluation of Using Large Language Models for Automated Unit Test
Generation,''
\textit{IEEE Trans. Softw. Eng.}, vol.~50, no.~1, pp.~85--105, 2024.

\bibitem{Yuan2023}
Z.~Yuan, Y.~Lou, M.~Liu, S.~Ding, K.~Wang, Y.~Chen, and X.~Peng,
``No More Manual Tests? Evaluating and Improving ChatGPT for Unit Test
Generation,''
\textit{arXiv preprint arXiv:2305.04207}, 2023.

\bibitem{Alshahwan2024}
N.~Alshahwan \textit{et al.},
``Automated Unit Test Improvement Using Large Language Models at Meta,''
in \textit{Proc. ACM FSE (Industry Track)}, 2024.

\bibitem{Yang2024}
L.~Yang \textit{et al.},
``On the Evaluation of Large Language Models in Unit Test Generation,''
\textit{arXiv preprint arXiv:2406.18181}, 2024.

\bibitem{Pereira2024}
A.~Pereira, B.~Lima, and J.~P. Faria,
``APITestGenie: Automated API Test Generation through Generative AI,''
\textit{arXiv preprint arXiv:2409.03838}, 2024.

\bibitem{Stennett2025}
T.~Stennett, M.~Kim, S.~Sinha, and A.~Orso,
``AutoRestTest: A Tool for Automated REST API Testing Using LLMs and MARL,''
in \textit{Proc. ICSE Companion}, 2025.

\bibitem{Kim2025}
M.~Kim \textit{et al.},
``LlamaRestTest: Effective REST API Testing with Small Language Models,''
\textit{arXiv preprint arXiv:2501.08598}, 2025.

\bibitem{Atlidakis2019}
V.~Atlidakis, P.~Godefroid, and M.~Polishchuk,
``RESTler: Stateful REST API Fuzzing,''
in \textit{Proc. IEEE/ACM ICSE}, pp.~748--758, 2019.

\bibitem{Arcuri2019}
A.~Arcuri,
``RESTful API Automated Test Case Generation with EvoMaster,''
\textit{ACM Trans. Softw. Eng. Methodol.}, vol.~28, no.~1, pp.~1--37, 2019.

\bibitem{Kim2022}
M.~Kim, Q.~Xin, S.~Sinha, and A.~Orso,
``Automated Test Generation for REST APIs: No Time to Rest Yet,''
in \textit{Proc. ACM ISSTA}, pp.~289--301, 2022.

\bibitem{Papadakis2019}
M.~Papadakis, M.~Kintis, J.~Zhang, Y.~Jia, Y.~Le Traon, and M.~Harman,
``Mutation Testing Advances: An Analysis and Survey,''
\textit{Adv. Comput.}, vol.~112, pp.~275--378, 2019.





\end{thebibliography}

\end{document}
%% ============================================================
%% END OF PAPER
%% ============================================================
